% -------------------------------------------------------------------------
% WBS ID: RES-ML-GEN
% Deliverable: Generalisation, Uncertainty and Interpretability
% Status: In progress
% Evidence status: Not assessed
% Review status: Not reviewed
% Last updated: 2026-06-30
% Dependencies: RES-ML-TRN, RES-ML-SUR, RES-DAT-CAS
% Downstream interfaces: RES-ML-INT, RES-ML-SPEC
% -------------------------------------------------------------------------

\section{RES-ML-GEN --- Generalisation, Uncertainty and Interpretability}
\label{sec:res-ml-gen}

\subsection{Purpose and Scope}
\label{sec:res-ml-gen-purpose}

This Deliverable restricts generalisation claims to the defined
event-conditioned domain. Cross-event performance is removed as an initial
acceptance requirement.

\subsection{Research Questions}
\label{sec:res-ml-gen-research-questions}

% TODO: State the questions that must be answered before the Deliverable
% can be accepted.

\subsection{Terminology and Definitions}
\label{sec:res-ml-gen-definitions}

% TODO: Define the technical terminology, symbols, quantities and
% classifications used in this Deliverable.

\subsection{Literature Evidence}
\label{sec:res-ml-gen-literature-evidence}

% TODO: Record evidence from peer-reviewed papers, authoritative datasets,
% standards, technical reports and primary documentation.
%
% Do not create a detached literature-summary list. Explain how each source
% supports, contradicts or limits a project decision.

\subsection{Governing Relations and Quantitative Definitions}
\label{sec:res-ml-gen-governing-relations}

% TODO: Record relevant equations, dimensionless groups, empirical models,
% extraction rules or quantitative definitions.
%
% Use align environments for displayed equations.
% Every displayed equation must have a unique \label{eq:...}.
% Do not place punctuation after displayed equations.

\subsection{Machine-Learning Role}
\label{sec:res-ml-gen-machine-learning-role}

% TODO: Complete this RES-ML Domain-specific subsection.

\subsection{Non-ML Baseline}
\label{sec:res-ml-gen-non-ml-baseline}

% TODO: Complete this RES-ML Domain-specific subsection.

\subsection{Input and Output Representation}
\label{sec:res-ml-gen-input-output-representation}

% TODO: Complete this RES-ML Domain-specific subsection.

\subsection{Dataset Requirements}
\label{sec:res-ml-gen-dataset-requirements}

% TODO: Complete this RES-ML Domain-specific subsection.

\subsection{Model or Architecture Candidates}
\label{sec:res-ml-gen-model-architecture-candidates}

% TODO: Complete this RES-ML Domain-specific subsection.

\subsection{Training and Validation Method}
\label{sec:res-ml-gen-training-validation-method}

% TODO: Complete this RES-ML Domain-specific subsection.

\subsection{Evaluation Metrics}
\label{sec:res-ml-gen-evaluation-metrics}

% TODO: Complete this RES-ML Domain-specific subsection.

\subsection{Generalisation and Uncertainty}
\label{sec:res-ml-gen-generalisation-uncertainty}

The methodology shall define interpolation limits, extrapolation detection,
held-out intervention tests, held-out severity tests, uncertainty calibration,
out-of-distribution rejection, feature sensitivity, engineering
interpretability, permitted-use boundaries, and conditions requiring a new FSI
simulation.

\subsection{Simulation and Optimisation Integration}
\label{sec:res-ml-gen-simulation-optimisation-integration}

% TODO: Complete this RES-ML Domain-specific subsection.

\subsection{Candidate Approaches}
\label{sec:res-ml-gen-candidate-approaches}

% TODO: Define the credible candidate approaches considered.

\subsection{Critical Comparison}
\label{sec:res-ml-gen-critical-comparison}

% TODO: Compare candidates using physically and project-relevant criteria.
% Do not select an approach solely on popularity or implementation ease.

\subsection{Assumptions and Applicability}
\label{sec:res-ml-gen-assumptions}

% TODO: State assumptions and the conditions under which they are valid.

\subsection{Limitations and Failure Conditions}
\label{sec:res-ml-gen-limitations}

The surrogate ceases to be applicable outside the validated event envelope,
outside represented defence and reinforcement types, where feature provenance is
insufficient, where predictive uncertainty is excessive, or where the
engineering consequence of an incorrect prediction requires direct simulation.

\subsection{Project-Specific Conclusions}
\label{sec:res-ml-gen-conclusions}

The initial project may claim only within-domain event-conditioned performance.
Cross-event generalisation remains future work.

\subsection{Selected Approach and Rationale}
\label{sec:res-ml-gen-selected-approach}

% TODO: Record the selected approach only after sufficient evidence exists.
% State the alternatives rejected and the reasons for rejection.

\subsection{Unresolved Decisions}
\label{sec:res-ml-gen-unresolved-decisions}

% TODO: Record unresolved questions, required evidence and decision owners.

\subsection{Requirements and Handoffs}
\label{sec:res-ml-gen-requirements}

% TODO: State requirements passed to other Research Domains, Software,
% verification activities or later execution work.

\subsection{Deliverable Completion Record}
\label{sec:res-ml-gen-completion-record}

% TODO: Record whether the Deliverable is:
%
% - Not started
% - In progress
% - Draft complete
% - Reviewed
% - Accepted
%
% Acceptance requires:
%
% - the research questions to be answered;
% - claims to be supported by appropriate evidence;
% - assumptions and limitations to be explicit;
% - the project-specific conclusion to be stated;
% - downstream requirements to be traceable;
% - unresolved decisions to be closed or formally deferred.
